% STATUS: written 20 October, expanded 26 October 2026.
%
% HOW FAR EACH SOURCE WAS READ. lion2024 and geotp2025 were read from their
% PDFs. lu2022rlrp was read in full for a separate course presentation in
% early November 2026; its internal details below rest on that reading. EVERY OTHER CITATION IN THIS SECTION WAS READ AT ABSTRACT LEVEL ONLY,
% through the Consensus index, and verified against Crossref. docs/literature-log.md
% records the searches, what was blocked, and the two entries Crossref does not
% index. kostler2023fluidity, vanbenschoten2022multiregion, eldeeb2022neuroshard
% and costafilho2020adaptive were added on 16 November 2026 from abstracts only
% (Consensus), with metadata copied from Crossref the same day.
% annamalai2018akkio was added the same day from its USENIX abstract page;
% USENIX papers carry no Crossref DOI.
%
% That limit governs what may be written here. A sentence summarizing a paper's
% abstract is supportable. A sentence about its internals, its baselines or its
% failure modes is not, unless the paper is opened first. Do not add one.

\section{Related Work}

Three bodies of work meet in this problem: the classical treatment of
fragmentation and allocation as an optimization to be solved, recent systems
that adapt placement to a moving workload, and the use of learning to replace
hand-written database administration.

\subsection{Fragmentation and allocation as an optimization problem}

The question of how to divide a database and where to put the pieces is as old
as distributed databases themselves, and the classical treatment is to state it
as an optimization problem and solve it once. A cost model scores a candidate
arrangement against an expected workload, and a solver searches for the
arrangement that minimizes it. The result is then installed and the database
runs with it.

The literature divides the problem in two. Fragmentation asks how to cut the
data; allocation asks where to put the pieces. Ceri et al.~\cite{ceri1982horizontal}
treat horizontal fragmentation, dividing a relation by rows, and Navathe et
al.~\cite{navathe1984vertical} treat vertical fragmentation, dividing it by
columns so that a fragment matches the attributes a transaction actually uses.
Both optimize against an expected access pattern supplied in advance. On the
allocation side, Apers~\cite{apers1988allocation} formalizes the placement of
fragments across network sites as a cost minimization and analyzes its
complexity, Sacc{\`a} and Wiederhold~\cite{sacca1985clustering} give a greedy
construction for a cluster of processors, and March and
Rho~\cite{march1995allocating} extend the objective to place operations as well
as data, under capacity constraints. The standard synthesis of this material is
the textbook treatment of {\"O}zsu and Valduriez~\cite{ozsu2011principles}, and
Nashat and Amer~\cite{nashat2018taxonomy} give a taxonomy of the techniques that
followed. The formulation remains current: Amer et
al.~\cite{amer2017optimized} solve fragmentation and allocation together as one
heuristic optimization, and Niu et al.~\cite{niu2026allocation} state the
allocation problem as a balanced cut.

A second tradition avoids the optimization entirely. Consistent
hashing~\cite{karger1997consistent} assigns a key to a node through a function
of the key alone, so that adding or removing a node moves only a small share of
the data and no central directory is needed. CRUSH~\cite{weil2006crush} extends
that to replicated data under user-defined failure domains, and a recent survey
compares the family~\cite{coluzzi2023hashing}. This is what the static baseline
in this paper is, and it is what most production systems use. It is also the
root of the problem studied here: a function of the key cannot observe where a
request came from.

A solved-once answer degrades as the real workload departs from the one it
was solved against: tolerable under slow drift, not when the workload's
geographic center moves daily. The systems in the next group re-solve the
problem as the workload changes.

\subsection{Adapting placement to a moving workload}

A line of work makes the partitioning decision from the observed workload
rather than from a design-time estimate. Schism~\cite{curino2010schism} builds a
graph in which tuples accessed by the same transaction are joined by an edge,
then partitions that graph to minimize distributed transactions whilst keeping
partitions balanced. Pavlo et al.~\cite{pavlo2012skew} search the design space
with an analytical cost model that accounts for skew as well as coordination
cost. E-Store~\cite{taft2014estore} addresses the case where a small number of
tuples carry most of the traffic, placing hot ranges individually whilst
distributing cold data in large chunks. Clay~\cite{serafini2016clay} removes the
restriction to hierarchical schemas, migrating blocks built around hot tuples
together with the tuples co-accessed with them. Kamal et
al.~\cite{kamal2016repartitioning} make the repartitioning incremental,
recutting only the sub-graph that distributed transactions touch so that the
scheme is refreshed as the workload drifts rather than recomputed.
Promi~\cite{ding2025promi} addresses the migration itself rather than the
decision, moving data at a finer granularity so that rebalancing does not stop
transactions.

A related line treats placement and replication as one problem rather than two
stages. Atrey et al.~\cite{atrey2019unifying} make that argument directly for
geo-distributed clouds, solving placement and replication jointly by overlapping
correlation clustering, and Adnan et al.~\cite{adnan2025network} optimize
placement against network cost in the same setting. Both reach the joint framing
this work uses, and both reach it with a solver rather than a learned policy.
Akkio~\cite{annamalai2018akkio}, in production at Facebook, migrates small
units of data between datacenters to cut access latency and cross-datacenter
traffic, and SLOG~\cite{ren2019slog} gives each granule of data a home region
by where it is accessed. Fluidity~\cite{kostler2023fluidity} moves the replicas of a replicated state
machine at run time towards the locations its client requests come from, and
CockroachDB's multi-region abstractions~\cite{vanbenschoten2022multiregion}
let a developer declare data locality in SQL, which the database maps to
replica placement and partitioning. In all four, geography drives placement,
and none is presented as a learned policy.

These systems share the premise of this work, that placement should follow the
workload, and they share a limitation. Each re-solves an optimization problem
when a trigger fires, using a cost model written by hand. None improves from the
outcomes of its own past decisions, and none of their cost models carries a term
for geographic distance, because they were designed for clusters in one data
center where the distance between any two nodes is much the same.

Lion~\cite{lion2024} is the closest system in intent. It arranges replicas so
that some single node holds every partition a transaction touches, letting it
commit locally: a graph analyzer groups partitions accessed together, a cost
model assigns the groups to nodes under a load-balance constraint, and an LSTM
forecasts the workload so replicas are provisioned before a shift. Lion reports
up to 2.7 times the throughput of the approaches it is compared against. It
shows that adapting to a moving workload pays, and that migration cost and
load balance belong in the objective. It leaves three openings: a hand-written
cost model decides, so it does not learn from its own outcomes; that model is
in bytes, with no distance term; and partitions are never split.

GeoTP~\cite{geotp2025} supplies the geography Lion lacks. For middleware over
data sources in several regions, a decentralized prepare phase cuts the round
trips of a distributed transaction and a latency-aware scheduler shortens the
time locks are held. Its cost accounting is the model for the one used here,
but it optimizes execution against a given placement and never moves data.
Lion adapts placement without seeing geography; GeoTP sees geography without
adapting placement; this work sits between them.

\subsection{Learning to administer a database}

A separate line of work replaces hand-written database administration with
learned policies. In configuration tuning, CDBTune~\cite{zhang2019cdbtune} treats
the setting of database knobs as a reinforcement learning problem and learns
from measured performance rather than from a regression over collected samples,
and QTune~\cite{li2019qtune} extends that to tune per query rather than per
workload; the survey of Zhao et al.~\cite{zhao2023knob} covers the area. In query
optimization, Neo~\cite{marcus2019neo} learns to produce execution plans and
improves from the plans it has run. UDO~\cite{wang2021udo} spans both, tuning
index selection and system parameters together and distinguishing parameters
that are expensive to change from those that are cheap, which is the same
distinction this work draws between moving a lease and moving a replica.

Placement has received less attention within this family than tuning or
optimization. RLRP~\cite{lu2022rlrp} is the nearest in method: a Deep Q-Network
places the replicas of each virtual node in an object store, a second agent
migrates replicas when a node joins, and the whole is validated on a real
system, Ceph. Its reward is load fairness, a placement is decided once per
virtual node, data moved is not charged, and request origin plays no part.
Two of its devices recur here: taking the next-highest action when the best is
invalid, which is masking at selection time, and reusing trained weights with
zero-initialized inputs when nodes are added, the natural route past this
work's fixed region count. Neuroshard~\cite{eldeeb2022neuroshard} learns shard assignments from past
queries with deep reinforcement learning, optimizing several objectives
together, fanout and load balance among them, and Costa Filho et al.~\cite{costafilho2020adaptive} change replica
placement in a key-value store incrementally with deep reinforcement learning
to balance load. Sibyl~\cite{singh2022sibyl} places data across
hybrid storage devices with online reinforcement learning, and
Harmonia~\cite{nadig2026harmonia} extends that to coordinate placement and
migration as two agents. Zheng et al.~\cite{zheng2024replica} place replicas in
an edge--cloud setting where geography does enter the objective, alongside
reliability and load, and Liao et al.~\cite{liao2019deep} learn placement
against a latency objective inside a data center. Closest of all in problem
statement, Mahmoud et al.~\cite{mahmoud2025rlfawm} formulate fragment allocation
and replication as a Markov decision process and solve it with tabular
Q-learning, driven by read, write and co-access matrices.

These differ from the present work in one or both of two ways. Most place data
across \emph{devices or tiers} that differ in speed, or across edge nodes that
differ in capacity, rather than across regions that differ in distance from the
traffic. And most hold the unit of placement fixed: a fragment, an object or a
block is placed, never split. Neuroshard learns the shard assignment itself,
with fanout and load balance as its stated objectives.
Mahmoud et al.~\cite{mahmoud2025rlfawm} come closest on the first axis and
still treat fragments as given, and a tabular
Q-function does not extend to the state space a six-region cluster with twenty
shards produces.

\subsection{Positioning}

This work gives one learned policy both fragmentation and geographic
placement actions, for transactional workloads. The workload-driven partitioners re-solve a
hand-written cost model and do not learn from their own outcomes. Lion learns a
forecast but not a policy, and its cost model is denominated in bytes rather
than in distance. GeoTP optimizes execution against a fixed placement. The
learned administration systems address knobs, plans and indexes. Akkio, SLOG, Fluidity and CockroachDB's multi-region design see geography but are
not presented as learned policies.
The learned placement systems address storage tiers, edge nodes and, in RLRP,
the even spread of load within one cluster, not a transactional workload whose
center of gravity moves, and only Neuroshard learns the shard assignment
itself.

Stating the contribution precisely: the agent chooses, in one action space,
where a shard's primary copy sits, where its replicas sit, and where the
boundary between two shards falls, and it is rewarded on a cost that combines
transaction latency, the volume of data moved, and load imbalance. That
splitting is what earns its result is not shown: splits are at the midpoint,
irreversible and capped at 20 shards, and no agent with splitting disabled was
trained for comparison.

It is not claimed to be new: searches of Semantic Scholar, OpenAlex and arXiv
found no paper that learns fragmentation and geographic placement jointly for
transactional workloads, but IEEE Xplore and the ACM Digital Library were
reached only through them. And on latency alone the agent does not beat the
heuristic it is meant to replace. What this paper establishes is the problem,
the apparatus, the baselines, and when a learned policy has room to win.
